%!TEX TS-program = xelatex
%!TEX encoding = UTF-8 Unicode

%===============================================================================
% FINAL SYSTEM-LOCKED RESUME: GOOGLE FDE OPTIMIZED
%===============================================================================

\documentclass[11pt, a4paper]{awesome-cv}

\usepackage{microtype}
\usepackage{enumitem}
\setlist[itemize]{topsep=1pt,itemsep=1pt,parsep=0pt,partopsep=0pt}

\geometry{left=1.3cm, top=.7cm, right=1.3cm, bottom=1.3cm, footskip=0cm}

\fontdir[fonts/]
\colorlet{awesome}{awesome-darknight}
\setbool{acvSectionColorHighlight}{false} 
\renewcommand{\acvHeaderSocialSep}{\quad|\quad}
\usepackage{hyperref}

%-------------------------------------------------------------------------------
% HEADER
%-------------------------------------------------------------------------------
\name{}{\href{udaylunawat.github.io}{Uday Lunawat}}
\position{\textbf{Senior Applied AI \& Agent Systems Engineer | RAG, Multi-Agent Systems, LLMOps}}
\mobile{\href{https://wa.me/917020901969}{(+91) 7020901969}}
\email{udaylunawat@gmail.com}
\homepage{udaylunawat.github.io}
\github{udaylunawat}
\linkedin{uday-lunawat}
\twitter{udaylunawat}

\begin{document}
\makecvheader[C]

%===============================================================================
% SUMMARY
%===============================================================================
\vspace{-1.2em}
\cvsection{Summary}

\begin{cvparagraph}
Senior Applied AI \& Agent Systems Engineer with \textbf{7 years of experience} building \textbf{production-grade systems around pretrained models (LLMs \& CV)}. Expertise in \textbf{Retrieval-Augmented Generation (RAG)}, prompt engineering, fine-tuning workflows, and \textbf{multi-agent orchestration} \\ using LangGraph.

Proven track record of delivering \textbf{customer-facing AI solutions} by integrating models with tools, memory, and APIs. Led development of an enterprise AI platform supporting \textbf{400+ production agents and 17,000+ deployments}, with strong focus on \textbf{evaluation and reliability}.
\end{cvparagraph}

\vspace{-1.3em}

%===============================================================================
% WORK EXPERIENCE
%===============================================================================
\cvsection{Work Experience}
\vspace{-0.5em}
\begin{cventries}

\cventry
{\underline{Tech}: Python, GCP, LangGraph, FastAPI, Jenkins, OIDC, SQLModel, Docker, Alembic}
{Senior AI Platform Engineer (GenAI, LLMOps, CI/CD) – \underline{\href{https://fractal.ai/}{Fractal Analytics}}}
{\textbf{Nov 2024 – Present}}
{Pune}
{
\begin{cvitems}
    \item {\textbf{Customer-Facing AI Systems:} Delivered agent-based AI solutions for enterprise stakeholders, integrating LLMs with tools, APIs, and \\ business workflows to solve real-world telecom use cases.}
    \item {\textbf{Pretrained Model Systems:} Built production systems around LLMs (OpenAI, Claude, Gemini) using prompt engineering, structured outputs, and tool augmentation.}
    \item {\textbf{Multi-Agent Systems:} Built stateful orchestration workflows (\textbf{LangGraph}) using ReAct and hierarchical delegation \\ for complex multi-step reasoning tasks.}
    \item {\textbf{Evaluation \& Observability:} Built evaluation pipelines (Langfuse, semantic validation) with LLM-native metrics to guide optimization \\ and system improvements.}
    \item {\textbf{Scale:} Platform supports \textbf{400+ production agents} and \textbf{17,000+ deployments}, reducing onboarding time to \textbf{<90 seconds}.}
    \item {\textbf{LLM FinOps:} Implemented cost-aware execution (token budgeting, adaptive context while maintaining accuracy SLAs.}
    \item {\textbf{Reliability Engineering:} Designed fault-tolerant systems (retries, fallbacks, partial recovery) improving success rate of multi-step workflows.}
    \item {\textbf{Security \& Governance:} Enforced Zero-Trust architecture (OIDC, RBAC, secrets isolation) for multi-tenant deployments.}
    \item {\textbf{Leadership:} Led a team of 7 engineers; drove architecture, reviews, and enterprise adoption.}
\end{cvitems}
}

\vspace{1.0em}

\cventry
{\underline{Tech}: Python, GCP, MLflow, Seldon Core, TensorFlow, Label Studio, Flask}
{Senior Machine Learning Engineer – \underline{\href{https://www.hcltech.com/}{HCL}}}
{\textbf{Apr 2024 – Jun 2024}}
{Bengaluru}
{
\begin{cvitems}
  \item {Architected \textbf{human-in-the-loop (HITL)} evaluation frameworks integrating feedback loops into model validation.}
  \item {Developed scalable annotation pipelines using Label Studio and microservices.}
  \item {Deployed ML inference systems with Docker and Cloud Run ensuring reliability and observability.}
\end{cvitems}
}

\vspace{1.0em}

\cventry
{\underline{Tech}: Python, AWS, MLflow, IaC, nbdev, Quantization}
{Machine Learning Engineer – \underline{Chugani}}
{\textbf{Feb 2022 – May 2023}}
{Bengaluru}
{
\begin{cvitems}
  \item {Built real-time liveness detection system achieving \textbf{150ms latency} via model optimization and quantization.}
  \item {Implemented MLflow for experiment tracking and modernized ML infrastructure using IaC.}
  \item {Refactored legacy ML pipelines into modular, documented Python packages using \textbf{nbdev}, improving maintainability, reuse, and collaboration.}
\end{cvitems}
}

\vspace{1.0em}

\cventry
{\underline{Tech}: Python, Airflow, AWS, ETL, Tableau}
{Machine Learning Engineer – \underline{Yang}}
{\textbf{Apr 2019 – Jan 2022}}
{Bengaluru}
{
\begin{cvitems}
  \item {Built automated AI audit system processing real-time factory data streams for quality control.}
  \item {Delivered dashboards across 20+ vendor sites saving \textbf{200+ hours/week}.}
\end{cvitems}
}

\vspace{1.0em}

\cventry{} {Systems Engineer Intern –\underline{\href{https://www.infosys.com/}{Infosys}}} {\textbf{Oct 2018 – Mar 2019}} {} {}
\vspace{-2.0em}
\cventry{} {Analytics Intern – \underline{\href{https://www.tcs.com/}{TCS}}} {\textbf{Jan 2018 – July 2018}} {} {}

\end{cventries}

\vspace{-3.5em}

%===============================================================================
% SKILLS
%===============================================================================
%-------------------------------------------------------------------------------
%	SKILLS
%-------------------------------------------------------------------------------
\cvsection{Skills}

\newenvironment{mylist}[1]
  {\begin{description}[labelwidth=#1,align=left,labelsep=0.6em,itemsep=-0.2em]}
  {\end{description}}

\begin{cventries}
\small
\vspace{-2.9em}

\cventry
    {}
    {}
    {}
    {}
    {
      \begin{mylist}{3.2cm}
        \item[Languages :] \textbf{Python}, SQL, Bash, Groovy
        \item[LLM \& GenAI :] Hugging Face Transformers, RAG, Prompt Engineering, Fine-tuning (LoRA/PEFT), LangGraph, LangChain
        \item[Agent Systems :] Multi-Agent Systems, ReAct, Self-Reflection, Tool Augmentation, MCP, Workflow Orchestration
        \item[Cloud \& Systems :] AWS, GCP, Docker, Kubernetes, FastAPI, Microservices, Terraform, CI/CD
        \item[MLOps \& Serving :] MLflow, Seldon Core, Model Serving, Experiment Tracking, Feature Pipelines
        \item[Evaluation \& FinOps :] LLM Metrics, Tracing (Langfuse), Cost Optimization
        \item[Data \& ML :] Pandas, NumPy, TensorFlow, PyTorch
        \item[AI Tooling :] OpenAI, Gemini, Claude APIs, GitHub Copilot, Claude Code
        \item[Certifications :] \href{https://verify.skilljar.com/c/x4j737hk74bn}{\textbf{Anthropic Claude Certified Architect}}, \href{https://github.com/udaylunawat/Data-Science-Projects}{AppliedAI Course}, \href{https://drive.google.com/file/d/1jT354cGLbs9ZUIO4jf71Qiy5wZDAVQKI/view?usp=sharing}{Pyspark and Data Engineering}
      \end{mylist}
    }

\end{cventries}
\vspace{-1.5em}

%===============================================================================
% EDUCATION
%===============================================================================
\cvsection{Education}
\vspace{-0.4em}

\begin{cventries}
\cventry{} {B.Tech. in Information Technology RCOEM, Nagpur} {\textbf{Apr 2014 – June 2018}} {} {}
\end{cventries}
\vspace{-3.6em}
\begin{cventries}
\cventry{} {Diploma in Computer Engineering} {\textbf{Apr 2011 – June 2014}} {} {}
\end{cventries}

\end{document}